\section{Conclusion}
\label{sec:conclusion}
% Our experiments show that the RL + LLM method outperforms both the Random + LLM and LLM baselines, achieving significantly lower solving times and higher solution rates, particularly for more complex constraint types. The analysis of time distribution and constraint breakdown reveals that the combination of reinforcement learning and large language models is particularly effective in guiding the search process and avoiding suboptimal variable assignments. The method also scales well to larger datasets, as demonstrated by its application to the SMT_COMP dataset, and shows promising results for the more challenging \QFNIA{} benchmarks.

% In summary, the proposed method not only improves the efficiency and accuracy of constraint solving but also demonstrates strong scalability, making it a robust solution for a wide range of SMT problems.
% \cite{even2020closer}

This work presented \tool, a solver-agnostic framework that casts SMT constraint simplification as policy learning: an \RL{} agent selects variables, an \LLM{} proposes values, and the solver verifies soundness in the loop. Across heterogeneous backends, \tool consistently increases solved counts and reduces time (\eg, Z3: 72→94, +30.6\%; CVC5: 228→360, +57.9\%; BVParti: 2→19, +850\%; MathSAT: 29→100, +244.8\%; see Section~\ref{sec:evaluation}). Ablations confirm the synergy between policy guidance and semantic value proposals, and cross-solver results (RQ4) indicate generalizability rather than solver-specific overfitting.

Looking ahead toward practical neuro-symbolic co-pilots~\cite{neurosymbolic2021}, we will: (i) expand theory coverage to bit-vectors, arrays, and quantified theories with solver-invariant abstractions; (ii) tighten LLM–solver integration via proof objects, unsoundness rejection, and expert feedback; (iii) improve deployment with calibrated, budget/risk-aware routing (RQ5) and policy distillation/caching to cut latency; and (iv) broaden datasets and standardized benchmarks to stress-test robustness and transfer.

% Despite its strengths, our method incurs additional computational overhead. Each LLM call introduces latency and resource cost that must be offset by downstream solver speedup. Training an effective RL policy can also be time-consuming; poor initial policies may explore excessively large search spaces before converging. Additionally, incorrect LLM suggestions can exacerbate solving difficulty, leading to cases where RL/LLM guidance inadvertently lengthens the solving process. Thus, the method is most effective when constraint solving is a clear bottleneck, and verifying the end-to-end performance impact is crucial.
